% !TeX program = xelatex
% Copyright (C) 2026 bsxiaocai
% SPDX-License-Identifier: LPPL-1.3c
% Example under LPPL 1.3c. See LICENSE and NOTICE.md.
% Compile from the repository root.
\documentclass[aspectratio=169]{ctexbeamer}
\usetheme{HUAS}
\HUASsetup{style=warm,logo=true,footer=minimal,sectionpage=true,
  faculty={读书分享与研讨}}
\input{examples/local-assets.tex}
\title{在文本中寻找解释的依据}
\subtitle{读书分享与研讨示例}
\author{分享者}
\institute{湖南文理学院}
\date{\today}
\begin{document}
\HUASmaketitle
\begin{frame}{分享结构}
  \tableofcontents[hideallsubsections]
\end{frame}
\section{阅读与细读}
\begin{frame}{先读原文，再提出问题}
  \begin{HUASquotation}{陆游《游山西村》}
    山重水复疑无路，柳暗花明又一村。
  \end{HUASquotation}
  \medskip
  \begin{itemize}
    \item 哪些词语表现行进过程与视野变化？
    \item “疑”与“又”怎样影响读者的感受？
    \item 除了抽象道理，诗句还呈现了怎样的景物？
  \end{itemize}
\end{frame}
\begin{frame}{文本细读：记录与解释并列}
  \begin{columns}[T,onlytextwidth]
    \begin{column}{0.47\textwidth}
      \begin{block}{可直接观察的语词}
        \begin{itemize}
          \item 山重、水复
          \item 疑无路
          \item 柳暗、花明、又一村
        \end{itemize}
      \end{block}
    \end{column}
    \begin{column}{0.47\textwidth}
      \begin{exampleblock}{据此展开的解释}
        从行路受阻到视野开阔，景物与情绪同时发生变化。
        解释应联系前后语句，不把诗句压缩为单一道理。
      \end{exampleblock}
    \end{column}
  \end{columns}
\end{frame}
\section{交流与反思}
\begin{frame}{研讨：比较不同解释}
  \begin{enumerate}
    \item 每人提出一个判断，并展示对应的语词。
    \item 同伴提出另一种可能的理解。
    \item 回到原文，比较两种解释的支持程度。
  \end{enumerate}
  \begin{alertblock}{讨论边界}
    不同理解可以并存，但都需要说明依据。
  \end{alertblock}
\end{frame}
\HUASkeypoint{阅读后的反思}{保留自己的问题，\\也听见其他的解释。}
\HUASclosingframe
\end{document}
